\documentclass[11pt]{amsart}
\usepackage[margin=1.1in]{geometry}
\usepackage{amsmath,amssymb,amsthm,mathtools}
\usepackage[T1]{fontenc}
\usepackage{lmodern}
\usepackage{microtype}
\usepackage{enumitem}
\usepackage{booktabs}
\usepackage[dvipsnames]{xcolor}
\usepackage[colorlinks=true,linkcolor=blue!50!black,citecolor=blue!50!black,urlcolor=blue!50!black]{hyperref}

\newtheorem{theorem}{Theorem}[section]
\newtheorem{proposition}[theorem]{Proposition}
\newtheorem{lemma}[theorem]{Lemma}
\newtheorem{corollary}[theorem]{Corollary}
\theoremstyle{definition}
\newtheorem{definition}[theorem]{Definition}
\newtheorem{setup}[theorem]{Setup}
\theoremstyle{remark}
\newtheorem{remark}[theorem]{Remark}

\newcommand{\QQ}{\mathbb{Q}}
\newcommand{\ZZ}{\mathbb{Z}}
\newcommand{\RR}{\mathbb{R}}
\newcommand{\CC}{\mathbb{C}}
\newcommand{\cO}{\mathcal{O}}
\newcommand{\fa}{\mathfrak{a}}
\newcommand{\Av}{\widehat{A}}
\newcommand{\ch}{\operatorname{ch}}
\newcommand{\Ext}{\operatorname{Ext}}
\newcommand{\Hom}{\operatorname{Hom}}
\newcommand{\End}{\operatorname{End}}
\newcommand{\NS}{\operatorname{NS}}
\newcommand{\Tr}{\operatorname{Tr}}
\newcommand{\Nm}{\operatorname{Nm}}
\newcommand{\SL}{\mathrm{SL}}
\newcommand{\bbar}{\overline}
\newcommand{\lrc}{\lrcorner}

\title{Fourier--Mukai Exclusion of Exponential Characters at Quartic CM Fields}
\date{Working note, 29 September 2026}

\begin{document}

\begin{abstract}
We prove that on an abelian eightfold of Weil type for a quartic CM field, no
perfect complex whose Chern character is a Weil class plus an exponential in
the polarisation, with inverse exponent in an explicit lattice, meets the
corrected propagation criterion.
\end{abstract}

\maketitle

\section{What this note proves}\label{sec:status}

The manuscript \cite{BMB} reduces the propagation statement for the split
families of a quartic CM field to one object: a perfect complex $E$ on a
member $A$ with $\ch(E)=N\omega+P$ and $\dim\Ext^{2}(E,E)=r(\ch E)$, where
$\omega$ is a Weil class, $P$ a polynomial in the classes $\eta_{s}$, and
$r$ the rank of a contraction (\cite[Theorem~19.21(v)]{BMB}; the notation is
recalled in \S\ref{sec:setup}). The least value a complex meeting this can
have is $88$ \cite[Corollary~19.23]{BMB}, and after the Weil tori
\cite[Proposition~19.24, Corollary~19.26]{BMB} two shapes remain at $88$:
\begin{align*}
  &\textup{(Exp)}\quad \ch(E)=N\omega+c+\Tr\bigl(a\,(e^{\lambda\Theta}-1)\bigr),
    \qquad a,\lambda\in F_{0}^{\times},\\
  &\textup{(Lin)}\quad \ch(E)=N\omega+c+\eta_{s},
    \qquad s\in F_{0}^{\times}.
\end{align*}
The manuscript notes that an autoequivalence of $D^{b}(A)$ changes the
character. This note builds the autoequivalences that act on the two real
places separately and uses them.

\begin{enumerate}[label=(\arabic*),leftmargin=2.2em,itemsep=2pt]
\item \emph{A place-wise Mukai group} (\S\ref{sec:group}). Tensoring by line
  bundles $L_{b}$ with $c_{1}=\eta_{b}$, and the conjugates by the
  Fourier--Mukai transform of tensoring by line bundles $M_{\beta}$ on
  $\Av$, act on $H^{*}(A)$ through two commuting copies of
  $\mathfrak{sl}_{2}$, one at each real place of $F_{0}$, with parameters in
  two explicit lattices $\cO'$ and $\fa$ of $F_{0}$. The transforms by the
  $M_{\beta}$ fix every Weil class, and all of them preserve $\dim\Ext^{2}$
  and $r$.
\item \emph{Exclusion} (Theorem~\ref{thm:main}). If $\lambda^{-1}\in\fa$,
  equivalently if the inverse of the $\QQ$-homomorphism $\varphi_{\eta_{\lambda}}$
  is a homomorphism $\Av\to A$, the shape (Exp) is carried to a character that
  stays of Hodge type on every Weil torus, and the argument of
  \cite[Theorem~19.12]{BMB} excludes it. The same holds at rank $104$ with
  the mixed exponential $e^{\eta_{\lambda}}$ added. The exponential shapes
  fall into classes indexed by $\lambda^{-1}$ modulo $\fa$, and the theorem
  removes the class $0$.
\item \emph{The linear shape} (\S\ref{sec:linear}). No element of
  $\SL_{2}(\CC)\times\SL_{2}(\CC)$ carries (Lin) to a Weil tori character,
  because it corresponds to $x^{3}y$ and not to a sum of two fourth powers.
  Its rank can be shifted by $4\Tr(s\fa)$; its transform by the Weyl element
  is never a shifted sheaf; on a product member it needs at least four
  exterior products; and a direct sum reaching $88$ has one summand
  carrying the whole character.
\end{enumerate}

What is not proved: the shape (Exp) with $\lambda^{-1}\notin\fa$ (in
particular every $\lambda$ with $\lambda^{-1}\notin\cO_{F_{0}}$, since
$\fa\subseteq\cO_{F_{0}}$), and the shape (Lin). So the propagation statement
$\mathrm{P2}_{\mathrm{split}}$ is not closed, and nothing here bears on (F2)
or (F3$'$) (\S\ref{sec:remains}). Section~\ref{sec:prior} checks and folds in
two earlier working notes, and \S\ref{sec:computations} lists the
computations.

\section{Setup}\label{sec:setup}

\begin{setup}\label{setup:quartic}
We use the notation of \cite[\S19]{BMB}. $F$ is a quartic CM field with
totally real subfield $F_{0}$, real places $\tau_{1},\tau_{2}$ and nontrivial
automorphism $\iota$ of $F_{0}$; for $x\in F_{0}$ write $x_{t}=\tau_{t}(x)$.
$A$ is a member of an $(F,2,\delta)$-family: an abelian eightfold with an
order of $F$ in $\End(A)$ and a polarisation $L$ with class $\theta$. Over
$\RR$, $H^{1}(A,\RR)=H_{1}\oplus H_{2}$ with $F_{0}$ acting on $H_{t}$ through
$\tau_{t}$, and $H_{t}\otimes\CC=V_{\sigma}\oplus V_{\bbar\sigma}$ for the
two embeddings $\sigma,\bbar\sigma$ of $F$ over $\tau_{t}$; each $V_{\sigma}$
has dimension $4$, with $V_{\sigma}^{1,0}$ and $V_{\sigma}^{0,1}$ of
dimension $2$. Since $F_{0}$ is fixed by the Rosati involution,
$\theta=\theta_{1}+\theta_{2}$ with $\theta_{t}\in\bigwedge^{2}H_{t}$
symplectic on $H_{t}$, and $\theta_{t}$ pairs $V_{\sigma}$ with
$V_{\bbar\sigma}$. For $g\in F_{0}[[T]]$ put
$\Tr(g(\Theta))=g^{\tau_{1}}(\theta_{1})+g^{\tau_{2}}(\theta_{2})$, with
$g^{\tau_{t}}$ obtained by applying $\tau_{t}$ to the coefficients; so
$\eta_{s}=\Tr(s\Theta)=s_{1}\theta_{1}+s_{2}\theta_{2}$, and
$\Tr(g(\Theta))$ is a rational class when $g$ has coefficients in $F_{0}$.
The line $\bigwedge^{4}V_{\sigma}$ is spanned by $\alpha_{\sigma}$, and
$\alpha_{\Sigma}=\prod_{\sigma\in\Sigma}\alpha_{\sigma}$. A nonzero rational
Weil class is $\omega=\sum_{\sigma}u_{\sigma}\alpha_{\sigma}$ with every
$u_{\sigma}\neq0$; its place parts are
$\omega^{(t)}=u_{\sigma}\alpha_{\sigma}+u_{\bbar\sigma}\alpha_{\bbar\sigma}$.
\end{setup}

For $\gamma\in H^{*}(A,\CC)$, $r(\gamma)$ is the rank of
$HT^{2}(A)\to H^{*}(A,\CC)$, $\xi\mapsto\xi\lrc\gamma$, where
$HT^{2}(A)=\bigwedge^{2}H^{0,1}\oplus H^{0,1}\otimes T\oplus\bigwedge^{2}T$
and $T=H^{0}(T_{A})$. For a perfect complex $E$ with $\ch(E)=N\omega+P$,
$N\neq0$ and $P$ a rational polynomial in the $\eta_{s}$,
\cite[Theorem~19.21(v)]{BMB} gives
\begin{equation}\label{eq:criterion}
  \dim\Ext^{2}(E,E)\ \ge\ r(\ch E)\ \ge\ 80,
\end{equation}
and equality in the first inequality, which we call \emph{the criterion},
implies $\mathbf{W}(F,2,\delta)$: the Weil classes are algebraic on every
abelian variety of $(F,2,\delta)$-Weil type. The rank is computed in closed form in
\cite[Theorem~19.22]{BMB}; the pure shape $c+N\omega$, of rank $80$, is
excluded by \cite[Theorem~19.12]{BMB}, so the least rank of a complex meeting
the criterion is $88$, and at $88$ the two shapes of \S\ref{sec:status} are
all that remain \cite[Proposition~19.24, Corollary~19.26]{BMB}.

The Weil tori are the complex tori $S_{F}$ of \cite[Setup~19.9]{BMB}: the
complex structures on $H^{1}(A,\RR)$ commuting with $F$ with each
$V_{\sigma}^{1,0}$ of dimension $2$. The classes $\alpha_{\Sigma}$ are of Hodge
type on every member of $S_{F}$, and \cite[Theorem~19.12]{BMB} states:
\emph{if $\ch(E)$ stays of Hodge type on every member of $S_{F}$ and
$\ch_{2}(E)\neq0$, then $\dim\Ext^{2}(E,E)>r(\ch E)$.}

\section{The place-wise Mukai group}\label{sec:group}

\begin{definition}\label{def:ops}
For $t=1,2$ let $L_{t}$ be the product with $\theta_{t}$ on
$H^{*}(A,\RR)=\bigwedge^{*}H_{1}\otimes\bigwedge^{*}H_{2}$, let
$\pi_{t}\in\bigwedge^{2}H_{t}^{\vee}$ be the bivector inverse to
$\theta_{t}$, let $\Lambda_{t}=-\iota_{\pi_{t}}$ be minus the contraction with
it, normalised so that $\Lambda_{t}\theta_{t}=4$, and let $H_{t}$ act on
$\bigwedge^{k}H_{t}$ by $k-4$.
\end{definition}

\begin{lemma}\label{lem:sl2}
Each $(L_{t},H_{t},\Lambda_{t})$ is an $\mathfrak{sl}_{2}$-triple, and the
operators of the two places commute. On the span of the powers of
$\theta_{t}$ the assignment
$\theta_{t}^{k}/k!\mapsto\binom{4}{k}x^{4-k}y^{k}$ identifies the action with
that on binary quartics, $L_{t}=y\,\partial_{x}$ and
$\Lambda_{t}=x\,\partial_{y}$; in particular $e^{\lambda\theta_{t}}\mapsto
(x+\lambda y)^{4}$, and $1$ and $\theta_{t}^{4}/24$ go to $x^{4}$ and
$y^{4}$.
\end{lemma}

\begin{proof}
The first statement is the Lefschetz $\mathfrak{sl}_{2}$ of the symplectic
space $(H_{t},\theta_{t})$ of dimension $8$, acting on the factor
$\bigwedge^{*}H_{t}$; operators on different tensor factors commute, being
even. The span of the powers of $\theta_{t}$ is the irreducible
representation generated by $1$, of highest weight $4$, and
$\Lambda_{t}\theta_{t}^{k}/k!=(5-k)\,\theta_{t}^{k-1}/(k-1)!$ matches
$x\partial_{y}\binom{4}{k}x^{4-k}y^{k}=(5-k)\binom{4}{k-1}x^{5-k}y^{k-1}$.
Items (A) of \S\ref{sec:computations} check the relations.
\end{proof}

\begin{lemma}\label{lem:weil}
$\Lambda_{t}\alpha_{\sigma}=0$ for every $\sigma$ and $t$, and
$L_{t}\alpha_{\sigma}=0$ for $\sigma$ over $\tau_{t}$. So $\omega^{(t)}$ is
invariant under the $\mathfrak{sl}_{2}$ of its own place and is a lowest
weight vector, of weight $-4$, for the other.
\end{lemma}

\begin{proof}
For $\sigma$ over $\tau_{t}$, $V_{\sigma}$ is a Lagrangian subspace of
$(H_{t}\otimes\CC,\theta_{t})$, since $\theta_{t}$ pairs it with
$V_{\bbar\sigma}$ and it has half the dimension. Every term of $\theta_{t}$
has a factor in $V_{\sigma}$, which repeats one of $\alpha_{\sigma}$, so
$\theta_{t}\alpha_{\sigma}=0$; and $\pi_{t}$ contracts a factor from
$V_{\sigma}$ with one from $V_{\bbar\sigma}$, of which $\alpha_{\sigma}$ has
none. At the other place $\alpha_{\sigma}$ has degree $0$, hence weight $-4$,
and $\Lambda_{t'}$ lowers the degree. Item (B) checks this.
\end{proof}

\begin{definition}\label{def:lattices}
Let $\varphi_{L}\colon A\to\Av$ be the polarisation. Put
\[
  \cO'=\{b\in F_{0}:\varphi_{L}\circ b\in\Hom(A,\Av)\},\qquad
  \fa=\{\beta\in F_{0}:\beta\circ\varphi_{L}^{-1}\in\Hom(\Av,A)\}.
\]
Both are subgroups of $F_{0}$. If $e$ is the exponent of $\ker\varphi_{L}$
and $R=F_{0}\cap\End(A)$, then $R\subseteq\cO'\subseteq e^{-1}R$ and
$eR\subseteq\fa\subseteq R\subseteq\cO_{F_{0}}$; so both are lattices of rank
$2$, and both are $R$-modules. If $L$ is principal, $\fa=\cO'=R$.
\end{definition}

The inclusions: $e=\psi\circ\varphi_{L}$ for a homomorphism $\psi$, which is
$e\varphi_{L}^{-1}$; $\beta\circ\varphi_{L}^{-1}\in\Hom(\Av,A)$ gives
$\beta=(\beta\circ\varphi_{L}^{-1})\circ\varphi_{L}\in\End(A)$; and
$\varphi_{L}\circ b\in\Hom(A,\Av)$ gives $eb=\psi\circ(\varphi_{L}\circ b)
\in\End(A)$. Composing with elements of $R$ preserves both conditions.

\begin{proposition}\label{prop:autoeq}
\begin{enumerate}[label=\textup{(\roman*)},leftmargin=2.4em,itemsep=2pt]
\item For $b\in\cO'$ there is a line bundle $L_{b}$ on $A$ with
  $c_{1}(L_{b})=\eta_{b}$, and $T_{b}=L_{b}\otimes(-)$ acts on $H^{*}(A,\QQ)$
  by $\exp(b_{1}L_{1}+b_{2}L_{2})$.
\item For $\beta\in\fa$ there is a line bundle $M_{\beta}$ on $\Av$ with
  $\varphi_{M_{\beta}}=\beta\circ\varphi_{L}^{-1}$, and
  $S_{\beta}=\Phi_{\mathcal{P}}^{-1}\circ(M_{\beta}\otimes-)\circ
  \Phi_{\mathcal{P}}$, with $\Phi_{\mathcal{P}}$ the Fourier--Mukai transform
  of the Poincar\'e bundle, is an autoequivalence of $D^{b}(A)$ acting by
  $\exp\bigl(\epsilon(\beta_{1}\Lambda_{1}+\beta_{2}\Lambda_{2})\bigr)$ for
  a sign $\epsilon=\pm1$ that does not depend on $A$ or $\beta$.
\end{enumerate}
Write $G(A)$ for the group generated by the $T_{b}$ and the $S_{\beta}$.
\end{proposition}

\begin{proof}
On a complex abelian variety the N\'eron--Severi group is the group of
symmetric homomorphisms to the dual \cite[\S2.5]{BL04}. Elements of $F_{0}$
are fixed by the Rosati involution, so $\varphi_{L}\circ b$ and
$\beta\circ\varphi_{L}^{-1}$ are symmetric; when they are homomorphisms they
are $\varphi$ of line bundles. The class of $L_{b}$ is the form
$\theta(b\,\cdot,\cdot)$, which is $b_{t}\theta_{t}$ on $H_{t}$, so
$\ch(L_{b})=e^{\eta_{b}}$ and (i) follows, $L_{1}$ and $L_{2}$ commuting.

For (ii), $\Phi_{\mathcal{P}}$ is an equivalence by Mukai \cite{Muk81}, and
on cohomology it is $x\mapsto p_{*}(q^{*}x\cdot e^{\ell})$ with $\ell$ the
Poincar\'e class, the Todd classes being trivial \cite[Proposition~6.1]{BMB}.
In a basis $e_{i}$ of $H^{1}(A)$ with dual basis $f_{i}$ of $H^{1}(\Av)$,
$\ell=\sum e_{i}f_{i}$ and $f_{i}e^{\ell}=\partial_{e_{i}}e^{\ell}$; since
$p_{*}$ extracts the top coefficient in the $e_{i}$ and kills derivatives,
integration by parts gives $f_{i}\cdot\Phi(x)=\pm\Phi(\iota_{e_{i}^{\vee}}x)$,
and for a class $y\in\bigwedge^{2}H^{1}(\Av)$,
\[
  y\cdot\Phi_{\mathcal{P}}(x)=\Phi_{\mathcal{P}}(\iota_{y}x),\qquad\text{so}\qquad
  \Phi_{\mathcal{P}}^{-1}\bigl(e^{y}\,\Phi_{\mathcal{P}}(x)\bigr)=\exp(\iota_{y})\,x ,
\]
with $y$ read as a bivector on $H^{1}(A)$ through $\ell$ (item (K) checks the
identity for abelian varieties of dimension at most three, signs included).
The class $c_{1}(M_{\beta})$, read in this way, is the bivector of the
alternating form of $\beta\circ\varphi_{L}^{-1}$, which is
$\beta_{1}\pi_{1}+\beta_{2}\pi_{2}$ up to one universal sign, because
$\varphi_{L}^{-1}$ has bivector $\pi_{1}+\pi_{2}$ up to that sign and $F_{0}$
acts on $H_{t}$ through $\tau_{t}$. For $L$ principal and $\beta=1$ this is the
dual principal polarisation, and $\iota_{\pi}\theta=\pm8=\pm(4+4)$, which is
the normalisation of Definition~\ref{def:ops}.
\end{proof}

\begin{proposition}\label{prop:invariance}
For $g$ in the group generated by the $\exp(xL_{t})$ and $\exp(x\Lambda_{t})$,
$x\in\CC$, $r(g\gamma)=r(\gamma)$ for every $\gamma$. Hence for every
$\Phi\in G(A)$ and every perfect complex $E$,
\[
  \dim\Ext^{2}(\Phi E,\Phi E)=\dim\Ext^{2}(E,E),\qquad
  r(\ch\Phi E)=r(\ch E),
\]
and $E$ meets the criterion if and only if $\Phi(E)$ does.
\end{proposition}

\begin{proof}
Let $U\subset\End H^{*}(A,\CC)$ be the span of the products with elements of
$H^{0,1}$ and of the contractions with elements of $T=T^{1,0}$. The operators
$\xi\lrc(\cdot)$, $\xi\in HT^{2}(A)$, span exactly the products of two
elements of $U$. Now $[L_{t},\iota_{v}]=-(\iota_{v}\theta_{t})\wedge$, and
$\iota_{v}\theta_{t}\in H^{0,1}$ for $v\in T^{1,0}$ because $\theta_{t}$ is of
type $(1,1)$; $[L_{t},y\wedge]=0$; $[\Lambda_{t},y\wedge]$ is, up to sign, the
contraction with $\pi_{t}(y)\in T^{1,0}$ for $y\in H^{0,1}$, since $\pi_{t}$ is
of type $(1,1)$; and $[\Lambda_{t},\iota_{v}]=0$. So $\operatorname{ad}L_{t}$
and $\operatorname{ad}\Lambda_{t}$ preserve $U$, conjugation by $g$ preserves
$U$ and $U\cdot U$, and $\dim(U\cdot U)\,g\gamma=\dim g^{-1}(U\cdot U)g\,
\gamma=\dim(U\cdot U)\,\gamma$. An equivalence preserves $\Ext^{2}$, and it
acts on Chern characters by its cohomological action. Item (E) checks the
invariance on the fifteen values of \cite[Theorem~19.22]{BMB} and on random
words.
\end{proof}

\section{Exclusion of the exponential shapes}\label{sec:main}

\begin{lemma}\label{lem:cusp}
For $\lambda,\beta\in\CC$ and each $t$,
\[
  \exp(\beta\Lambda_{t})\,e^{\lambda\theta_{t}}=
  \begin{cases}
    (1+\lambda\beta)^{4}\,e^{\lambda'\theta_{t}},\quad
    \frac{1}{\lambda'}=\frac{1}{\lambda}+\beta, & 1+\lambda\beta\neq0,\\[2pt]
    \lambda^{4}\,\theta_{t}^{4}/24, & 1+\lambda\beta=0 .
  \end{cases}
\]
\end{lemma}

\begin{proof}
By Lemma~\ref{lem:sl2}, $\exp(\beta x\partial_{y})$ substitutes $y+\beta x$
for $y$, and $(x+\lambda(y+\beta x))^{4}=((1+\lambda\beta)x+\lambda y)^{4}$.
Item (C) checks both cases.
\end{proof}

\begin{theorem}\label{thm:exclusion}
Let $E$ be a perfect complex on $A$ and $\beta\in\fa$. If
$\exp(\beta_{1}\Lambda_{1}+\beta_{2}\Lambda_{2})\ch(E)$ stays of Hodge type on
every member of $S_{F}$ and has nonzero component in $H^{4}$, then
$\dim\Ext^{2}(E,E)>r(\ch E)$.
\end{theorem}

\begin{proof}
$\fa=-\fa$, so with $\epsilon$ as in Proposition~\ref{prop:autoeq} the
autoequivalence $\Phi=S_{\epsilon\beta}$ has $\ch(\Phi E)=
\exp(\beta_{1}\Lambda_{1}+\beta_{2}\Lambda_{2})\ch(E)$. By
\cite[Theorem~19.12]{BMB}, $\dim\Ext^{2}(\Phi E,\Phi E)>r(\ch\Phi E)$, and
both sides are those of $E$ by Proposition~\ref{prop:invariance}.
\end{proof}

\begin{theorem}\label{thm:main}
Let $A$ be a member of an $(F,2,\delta)$-family, $\omega$ a nonzero rational
Weil class, $N\neq0$, $c,b\in\QQ$, $a\in F_{0}$ and $\lambda\in F_{0}^{\times}$,
and let $E$ be a perfect complex on $A$ with
\[
  \ch(E)=N\omega+c+\Tr\bigl(a\,(e^{\lambda\Theta}-1)\bigr)
  +b\,(e^{\eta_{\lambda}}-1).
\]
If $\lambda^{-1}\in\fa$, equivalently if $\varphi_{\eta_{\lambda}}^{-1}
=\lambda^{-1}\circ\varphi_{L}^{-1}$ is a homomorphism $\Av\to A$, then $E$
does not meet the criterion: $\dim\Ext^{2}(E,E)>r(\ch E)$. Here
$r(\ch E)=88$ when $b=0\neq a$, and $r(\ch E)=104$ when $b\neq0$.
\end{theorem}

\begin{proof}
Take $\beta=-\lambda^{-1}\in\fa$. By Lemma~\ref{lem:weil} the operators
$\Lambda_{t}$ kill $\omega$ and $1$; by Lemma~\ref{lem:cusp} at both places,
the operators of the two places acting on separate tensor factors,
\[
  \exp(\beta_{1}\Lambda_{1}+\beta_{2}\Lambda_{2})\ch(E)
  =N\omega+(c-\Tr a-b)+\tfrac{1}{24}\Tr\bigl(a\lambda^{4}\Theta^{4}\bigr)
  +\tfrac{b}{576}\Nm(\lambda)^{4}\,\theta_{1}^{4}\theta_{2}^{4}.
\]
As in the proof of \cite[Proposition~19.24]{BMB}, $\theta_{t}^{4}$ is a
nonzero multiple of $\alpha_{\sigma}\alpha_{\bbar\sigma}$ for the embeddings
over $\tau_{t}$, so this class lies in the span of the $\alpha_{\Sigma}$, is
rational, stays of Hodge type on $S_{F}$, and has component $N\omega\neq0$ in
$H^{4}$. Theorem~\ref{thm:exclusion} applies. For the ranks, in the notation
of \cite[Theorem~19.22]{BMB} the Hankel matrices of $(a_{t}\lambda_{t}^{k})$
have rank $\rho_{t}=1$, so $R_{t}=8$, and $\mu=0$ exactly when $b=0$; the
formula gives $64+4+4+8+8=88$ and $88+16=104$. Items (D) and (F) check the
transform and the ranks.
\end{proof}

\begin{corollary}\label{cor:cusps}
Let $\beta\in\fa$ with $1+\lambda\beta\neq0$. Then $S_{\pm\beta}$ carries a
character of shape \textup{(Exp)} with parameters $(c,a,\lambda)$ to one with
\[
  a'=a(1+\lambda\beta)^{4},\qquad \frac{1}{\lambda'}=\frac{1}{\lambda}+\beta,
  \qquad c'=c+\Tr(a'-a).
\]
So the orbits of the shapes \textup{(Exp)} under the $S_{\beta}$ are indexed
by $\lambda^{-1}$ in $F_{0}/\fa$, an infinite group, and
Theorem~\ref{thm:main} removes exactly the class of $0$: a transform
$S_{\beta}$ puts \textup{(Exp)} in the span of the $\alpha_{\Sigma}$ only when
$1+\lambda\beta=0$.
\end{corollary}

\begin{proof}
Lemma~\ref{lem:cusp}. For $1+\lambda\beta\neq0$ the transform still has a
nonzero component $a'\lambda'\theta$ in the span of $\theta_{1},\theta_{2}$,
which is not in the span of the $\alpha_{\Sigma}$.
\end{proof}

\begin{corollary}\label{cor:principal}
If $L$ is principal and $\cO_{F_{0}}\subseteq\End(A)$, then $\fa=\cO_{F_{0}}$,
and Theorem~\ref{thm:main} excludes \textup{(Exp)} exactly for
$\lambda^{-1}\in\cO_{F_{0}}$: for instance $\lambda=1$, that is
$\ch(E)=N\omega+c+\Tr(a(e^{\Theta}-1))$, and $\lambda=1/m$, $m\in\ZZ\smallsetminus\{0\}$; not
$\lambda=2$. At any member, $\fa\subseteq\cO_{F_{0}}$, so no $\lambda$ with
$\lambda^{-1}\notin\cO_{F_{0}}$ is excluded by this argument.
\end{corollary}

\begin{remark}[Where the argument stops]\label{rem:swap}
The transforms in $G(A)$ act through the group $M$ of matrices
$\left(\begin{smallmatrix}u&b\\c&v\end{smallmatrix}\right)\in\SL_{2}(F_{0})$
with $u,v\in R$, $b\in\cO'$ and $c\in\fa$, embedded by
$(\tau_{1},\tau_{2})$; it is a group because $\cO'$ and $\fa$ are
$R$-modules and $\cO'\fa\subseteq R$. At each place the Weil part
of the other place, a lowest weight vector, pins the line of $x$, so a
transform that puts \textup{(Exp)} into the span of the $\alpha_{\Sigma}$
must send the pair of lines $\{x,\,x+\lambda y\}$ to $\{x,\,y\}$ at both
places. If it fixes the line of $x$, it is lower triangular with a unit on
the diagonal and an entry of $\fa$ below it, and it sends $x+\lambda y$ to the
line of $y$ exactly when $\lambda^{-1}\in\fa$: that is
Theorem~\ref{thm:main}. If it exchanges the two lines, it moves
$\omega^{(t)}\otimes x^{4}$ to $\omega^{(t)}\theta_{t'}^{4}/24$, so the
transformed character has no component in $H^{4}$ and a nonzero one in
$H^{12}$; \cite[Theorem~19.12]{BMB} does not apply, since the vanishing on a
very general Weil torus \cite[Proposition~19.11]{BMB} covers $\ch_{k}$ for
$1\le k\le3$ only. When $L$ is principal, the Weyl element lies in $M$, such an
exchange is the Weyl element times a transform of the first kind, and it
reaches no further $\lambda$. So Theorem~\ref{thm:main} is all this group
gives on \textup{(Exp)}, and $\lambda^{-1}\notin\fa$ needs another idea.
\end{remark}

\section{The linear shape}\label{sec:linear}

Throughout, $\gamma=N\omega+c+\eta_{s}$ with $N\neq0$, $c\in\QQ$ and
$s\in F_{0}^{\times}$, so $s_{1}s_{2}\neq0$.

\begin{proposition}\label{prop:nolift}
No element of $\SL_{2}(\CC)\times\SL_{2}(\CC)$, acting through the two
$\mathfrak{sl}_{2}$ of Lemma~\ref{lem:sl2}, carries $\gamma$ into the span of
the $\alpha_{\Sigma}$. In particular no $\Phi\in G(A)$ does, and
Theorem~\ref{thm:exclusion} does not apply to \textup{(Lin)}.
\end{proposition}

\begin{proof}
The group preserves the isotypic components of $H^{*}(A)$ for
$\mathfrak{sl}_{2}\times\mathfrak{sl}_{2}$. The component of
$\mathrm{Sym}^{4}\otimes\mathrm{Sym}^{4}$ is the span of the monomials
$\theta_{1}^{i}\theta_{2}^{j}$, because the only class of weight $(-4,-4)$ is
the constant; its intersection with the span of the $\alpha_{\Sigma}$ is
spanned by $x^{4}$ and $y^{4}$ at each place, that is by $1$,
$\theta_{1}^{4}$, $\theta_{2}^{4}$ and $\theta_{1}^{4}\theta_{2}^{4}$,
while the other $\alpha_{\Sigma}$ lie in other components
(Lemma~\ref{lem:weil}). The component of $\gamma$ there is
\[
  p=c\,x_{1}^{4}x_{2}^{4}+4s_{1}\,x_{1}^{3}y_{1}\,x_{2}^{4}
   +4s_{2}\,x_{1}^{4}\,x_{2}^{3}y_{2}.
\]
Read $p$ as a map from the dual of the quartics of the second place to those
of the first; its image contains $x_{1}^{3}y_{1}$ because $s_{1}s_{2}\neq0$.
If $(g_{1},g_{2})\,p$ lay in $\langle x^{4},y^{4}\rangle^{\otimes2}$, the
image would lie in $\langle m^{4},n^{4}\rangle$ with $m=g_{1}^{-1}x$,
$n=g_{1}^{-1}y$ independent. But a nonzero $um^{4}+vn^{4}$ is a fourth power
when $uv=0$ and has four distinct roots otherwise, its discriminant being
$256u^{3}v^{3}\det(m,n)^{12}$, while $x^{3}y$ has a triple root. Item (H)
checks the discriminant and the image.
\end{proof}

The exponential shape corresponds to the fourth power $(x+\lambda y)^{4}$,
which the group can move to $y^{4}$; the linear shape corresponds to its
derivative at $\lambda=0$, which it cannot.

\begin{proposition}\label{prop:rankshift}
For $\beta\in\fa$, $S_{\pm\beta}(E)$ has character
$N\omega+\bigl(c+4\Tr(s\beta)\bigr)+\eta_{s}$. So the rank of a complex of
shape \textup{(Lin)} meeting the criterion can be moved by any element of
$4\Tr(s\fa)\subseteq\ZZ$ without leaving the shape or the criterion; rank
$0$ is reached exactly when $c\in4\Tr(s\fa)$. The $\beta$ with
$\Tr(s\beta)=0$ form a subgroup of rank one of $\fa$, which fixes the
character.
\end{proposition}

\begin{proof}
$\exp(\beta\Lambda_{t})\theta_{t}=\theta_{t}+4\beta$ and $\Lambda_{t}$ kills
$1$ and $\omega$; Proposition~\ref{prop:invariance}. Item (G).
\end{proof}

\begin{proposition}\label{prop:notsheaf}
Let $L$ be principal, so that $1\in\cO'\cap\fa$, and let
$W=T_{1}\,S_{-\epsilon}\,T_{1}$, which acts by
$w=\exp(L)\exp(-\Lambda)\exp(L)$ at both places. For every $\beta\in\fa$
and every $E$ with $\ch(E)=\gamma$, the object $W S_{\beta}(E)$ is not
isomorphic to a shift of a coherent sheaf.
\end{proposition}

\begin{proof}
At each place $w$ sends $x\mapsto y$ and $y\mapsto-x$; so $w(1)$ is the
class $\theta_{1}^{4}\theta_{2}^{4}/576$ of a point, $w(\theta_{t})$ has
degree $14$, and $w(\omega)=\bigl(\omega^{(1)}\theta_{2}^{4}+
\theta_{1}^{4}\omega^{(2)}\bigr)/24$. By Proposition~\ref{prop:rankshift}
the character of $W S_{\beta}(E)$ therefore vanishes below degree $12$ and
has
\[
  \ch_{6}=\tfrac{N}{24}\bigl(\omega^{(1)}\theta_{2}^{4}+\theta_{1}^{4}\omega^{(2)}\bigr)
  \neq0,\qquad \int_{A}\ch_{6}\cdot\theta_{i}\theta_{j}=0\ \text{ for all }i,j,
\]
the integrals vanishing because $\theta_{t}\,\omega^{(t)}=0$
(Lemma~\ref{lem:weil}) and $\theta_{t}^{5}=0$. If $W S_{\beta}(E)\cong
\mathcal{F}[k]$ for a sheaf $\mathcal{F}$, then $\ch(\mathcal{F})$ vanishes
below degree $12$; the lowest nonzero Chern character of a sheaf is the
class of the cycle of its support in that codimension
\cite[\S15.1, Ex.~15.1.5 and Cor.~18.3.2]{Ful98}, which is effective, so $\mathcal{F}$ is
supported in dimension at most two and $\ch_{6}(\mathcal{F})$ is a nonzero
effective cycle, whose integral against $\theta^{2}$ is positive. This
contradicts the vanishing. Item (I) checks $w$ and the integrals.
\end{proof}

So a complex of shape (Lin) meeting the criterion cannot be produced as the
inverse transform, by $WS_{\beta}$, of a sheaf on a surface: the object must
be built as a genuine complex on that side as well.

\begin{proposition}\label{prop:product}
Suppose $A=B\times B'$ with $B,B'$ abelian fourfolds of $(F,1)$-Weil type and
the product structure. Then the component of $\ch_{2}(E)=N\omega$ in
$H^{2}(B)\otimes H^{2}(B')$ has tensor rank $4$. Hence if
$[E]=\sum_{i=1}^{k}\pm[E_{i}\boxtimes E'_{i}]$ in $K_{0}(A)_{\QQ}$, then
$k\ge4$; in particular $E$ is not an exterior product, nor a cone of two.
\end{proposition}

\begin{proof}
$\bigwedge^{4}V_{\sigma}=\bigwedge^{2}V_{\sigma}(B)\otimes
\bigwedge^{2}V_{\sigma}(B')$, so each $\alpha_{\sigma}$ is a pure tensor, and
the four lie in different eigenspaces of $F$ on both sides; the
$\theta_{t}$ have no component there. The component of
$\ch(E_{i}\boxtimes E'_{i})$ is $c_{1}(E_{i})\otimes c_{1}(E'_{i})$. Item
(J).
\end{proof}

\begin{proposition}\label{prop:sums}
Let $E=\bigoplus_{i}E_{i}$ meet the criterion, with each $\ch(E_{i})$ of the
form $N_{i}\omega+P_{i}$, $P_{i}$ a rational polynomial in the $\eta_{s}$.
Then each $E_{i}$ meets its own criterion, $\Ext^{2}(E_{i},E_{j})=0$ for
$i\neq j$, and $r(\ch E)=\sum_{i}r(\ch E_{i})$. If $r(\ch E)<108$, exactly one
summand has nonzero Chern character.
\end{proposition}

\begin{proof}
$\dim\Ext^{2}(E,E)\ge\sum_{i}\dim\Ext^{2}(E_{i},E_{i})\ge\sum_{i}
r(\ch E_{i})\ge r(\ch E)$, the last because a rank is subadditive; equality
throughout gives the first three statements. By
\cite[Theorem~19.21(iii)]{BMB}, $r(N_{i}\omega+P_{i})\ge80$ when
$N_{i}\neq0$. When $N_{i}=0\neq P_{i}$, $r(P_{i})\ge28$: the torus scaling
the classes of type $(1,0)$ at one place multiplies $\theta_{t}$ alone and
normalises $U$, so it keeps $r$, and degenerating $P_{i}$ to a vertex of its
Newton polygon, $r$ being lower semicontinuous, gives
$r(P_{i})\ge\min_{i,j}r(\theta_{1}^{i}\theta_{2}^{j})=28$ (item (L)).
Since $\sum N_{i}=N\neq0$, a second summand with nonzero character would give
$r\ge80+28$.
\end{proof}

\section{What remains}\label{sec:remains}

\begin{enumerate}[label=(\alph*),leftmargin=2.2em,itemsep=2pt]
\item \emph{$\mathrm{P2}_{\mathrm{split}}$ at a quartic field.} At rank $88$:
  the shape (Exp) with $\lambda^{-1}\notin\fa$, which includes every
  $\lambda$ with $\lambda^{-1}\notin\cO_{F_{0}}$, and the shape (Lin),
  subject to Propositions~\ref{prop:nolift}--\ref{prop:sums}. Above $88$,
  only the mixed exponential of rank $104$ with $\lambda^{-1}\in\fa$ is removed
  here. Within $G(A)$ the only further gain would come from the vanishing of
  $\ch_{6}$ on a very general Weil torus
  (Remark~\ref{rem:swap}) together with a transform exchanging the two lines,
  which exists only for non-principal $L$. What (Lin) needs is an object: a complex, not a
  shifted sheaf after $WS_{\beta}$, not a single summand of lower rank, and on
  a product member not fewer than four exterior products in $K$-theory.
\item \emph{(F2).} The mechanism of \S\ref{sec:group} needs a totally real
  field in $\End(A)$ to split the Lefschetz $\mathfrak{sl}_{2}$ into places. A
  Mumford fourfold has $\End=\ZZ$, its exceptional classes are rigid
  \cite[Theorem~21.21]{BMB}, and every object built from line bundles by
  cones, tensor products, duals, pull-backs, push-forwards and Fourier--Mukai
  functors with such kernels has a Chern character invariant under the
  Lefschetz group, which no exceptional class is
  \cite[Theorem~21.36]{BMB}. So no transform of the kind used here reaches
  (F2), and we have no construction for it.
\item \emph{(F3$'$).} Nothing here bears on it.
\end{enumerate}

\section{Earlier working notes}\label{sec:prior}

\subsection{The smooth-support note}\label{ssec:smooth}
A working note of the same date, \emph{Record of an attempt on the
smooth-support case at $n=4$}, concerns an imaginary quadratic field
$K=\QQ(\sqrt{-d})$ and a secant object $\mathcal{I}_{S}(3\Theta)$ with smooth
support on a principally polarised abelian fourfold $(X,\Theta)$, the shape
isolated in \cite[Corollary~18.68, Remark~18.69]{BMB} for Markman's
construction \cite{Mar25a,Mar25b}. That route is not needed by any minimal
sufficient set of the closure graph, since (P2) for the imaginary quadratic
fields follows from $\mathrm{P2}_{\mathrm{split}}$ by \cite[Proposition~20.13]{BMB}
and \cite[Theorem~20.58]{BMB}; it is recorded because what it proves is
correct, and it is now part of the manuscript. We checked its numerical
claims in exact arithmetic (the second program of \S\ref{sec:computations},
twenty-five checks) and read its proofs; the one claim not checked, that the
degree of the Gauss map divides $\gcd([S]^{2},e(S))$, is used neither here
nor in the manuscript. Write $u_{\Theta}=\sum_{k}(-d)^{k}\Theta^{2k}/(2k)!$ and
$v_{\Theta}=\sum_{k}(-d)^{k}\Theta^{2k+1}/(2k+1)!$, let
$\Lambda_{\Theta}\subset\QQ[\Theta]/(\Theta^{5})$ be the group generated by the
$e^{j\Theta}$, and for $\xi=\sum_{i}a_{i}\Theta^{i}/i!$ define the binomial
moments $m_{k}(\xi)$ by $a_{i}=\sum_{k}S(i,k)\,k!\,m_{k}$.

\begin{theorem}[{\cite[Lemma~18.56, Theorem~18.57]{BMB}}]\label{thm:lattice}
$\xi\in\Lambda_{\Theta}$ if and only if $m_{0},\dots,m_{4}\in\ZZ$, and
\[
  m\bigl(a\,u_{\Theta}+b\,v_{\Theta}\bigr)=\Bigl(a,\ b,\ -\tfrac{ad+b}{2},\
  \tfrac{3ad+2b-bd}{6},\ \tfrac{ad^{2}-11ad+6bd-6b}{24}\Bigr).
\]
In particular $u_{\Theta}+3v_{\Theta}\in\Lambda_{\Theta}$ if and only if
$d\equiv15$ or $23\pmod{24}$.
\end{theorem}

The proof is that $e^{j\Theta}\mapsto(1+y)^{j}$ identifies $\Lambda_{\Theta}$
with $\ZZ[y]/(y^{5})$ and the moments with the coefficients. It explains
\cite[Table~4]{BMB}: a split Hilbert--Burch resolution has its character in
$\Lambda_{\Theta}$, so only $d\equiv15,23\pmod{24}$ can occur, at any rank and
with any twists, which is what the search found below $24$. For the four
discriminants $d\in\{1,3,5,7\}$ of a smooth support
\cite[Corollary~18.41]{BMB} it excludes every resolution
$0\to E_{1}\to E_{0}\to\mathcal{I}_{S}\to0$ with $\ch E_{0},\ch E_{1}\in
\Lambda_{\Theta}$ \cite[Corollary~18.58]{BMB}, for instance by bundles built
from line bundles algebraically equivalent to multiples of $\Theta$ by sums,
tensor products, duals, exterior and symmetric powers, extensions, kernels of
surjections and cokernels of injections. These need not be projectively flat,
so for the smooth shape this goes beyond \cite[Theorems~18.63 and~18.66]{BMB}.
The original listed kernels and cokernels of arbitrary maps; every bundle is
the image of a map of constant rank between sums of line bundles, so those
can have any class, and the manuscript states the exact cases above.

The note also computes $m(\cO_{S})=(0,0,N,-3N,N(47-2N)/12+3N)$, whose last
entry is not an integer for $N\in\{5,6,7,8\}$, and reproduces
\cite[Table~3]{BMB} from Grothendieck--Riemann--Roch: $K_{S}^{2}=12N(36-N)$,
$e(S)=12N(36-3N)$, $\chi(\cO_{S})=4N(18-N)$.

\begin{lemma}[{\cite[Lemma~18.43]{BMB}}]\label{lem:inj}
Let $X$ be a simple abelian fourfold and $S\subset X$ a smooth surface with
$e(S)\neq[S]^{2}$. Then $H^{2}(\cO_{X})\to H^{2}(\cO_{S})$ is injective. For a
smooth secant support $[S]^{2}=24N^{2}\neq12N(36-3N)=e(S)$ for every integer
$N\ge1$, so on a simple $X$ the injectivity required in
\cite[Corollary~18.68(ii)]{BMB} holds automatically.
\end{lemma}

The proof: by conjugation the kernel is that of the restriction of invariant
$2$-forms; a form of rank two vanishing on $S$ gives, by Castelnuovo--de
Franchis, a fibration whose fibres generate an abelian subvariety of
dimension at most two, against simplicity; a form of rank four vanishing on
$S$ makes $S$ Lagrangian, so $N_{S/X}\cong\Omega^{1}_{S}$ and
$[S]^{2}=c_{2}(N_{S/X})=e(S)$. Simplicity is needed only in the rank-two
case, and the manuscript states it wherever the injectivity is called
automatic.

The note further shows that $\delta=K_{S}-4\lambda$, $\lambda=\Theta|_{S}$, is
a nonzero algebraic class with $\iota_{*}\delta=0$ and
$\delta^{2}=-12N(N-4)$ \cite[Remark~18.44]{BMB}, and, through Bloch's
semiregularity map \cite{Blo72}, that $S$ extends to first order along every
polarised deformation of $X$ if and only if that map has rank $6$, with
$h^{1}(N_{S/X})\ge88,296,568,904$ for $N=5,6,7,8$
\cite[Proposition~18.46]{BMB}. In the lower bound $h^{0}(N_{S/X})\ge4$ a
translation field tangent to $S$ everywhere would be a nowhere vanishing
vector field on $S$, against $e(S)=12N(36-3N)\ne0$; this is the argument the
manuscript uses, and it needs neither simplicity nor that $S$ be of general
type. The existence of $S$ and the first-order condition stay open.

Its last section, on the Kuga--Satake class of the Mumford square, proves
that for real multiplication of odd degree every Hodge similarity is a
rational scalar; as that note says, this is already in the manuscript. Of the
works it cites there, those of Morrison \cite{Mor84}, Buskin \cite{Bus19} and
Varesco \cite{Var23} are in the manuscript's bibliography; the others are not
repeated here.

\subsection{The status note}\label{ssec:statusnote}
A second note searched, by simulated annealing, unions of coordinate abelian
surfaces in $E_{1}\times\dots\times E_{4}$ for configurations with the
numerical invariants of a secant object at $b=1$. It reports exact matches at
$d=1$ and none at $d=2,3$. Its own caveats stand: $E_{1}\times\dots\times
E_{4}$ is not a member of a Weil family with the required action, and the
search does not touch the second condition of \cite[Theorem~18.29]{BMB}. It is
a data point about the imaginary quadratic route, which, as above, the closure
graph does not need.

\section{Computations}\label{sec:computations}

Two programs in exact rational arithmetic accompany this note. They are
kept with its source in the repository \url{https://github.com/creelie/H8},
in the directory \texttt{notes/fourier\_mukai}. The first reuses the exterior-algebra
model of \cite[item (LXV)]{BMB} and checks: (A) the two
$\mathfrak{sl}_{2}$-triples and their commutation; (B)
Lemma~\ref{lem:weil}; (C) Lemma~\ref{lem:cusp}; (D) the transform in the
proof of Theorem~\ref{thm:main} on five characters, and that it lands in the
span of the $\alpha_{\Sigma}$; (E) Proposition~\ref{prop:invariance} on the
fifteen values of \cite[Theorem~19.22]{BMB} and on random words; (F) the ranks
$88$ and $104$; (G) Proposition~\ref{prop:rankshift}; (H) the discriminant and
the image in Proposition~\ref{prop:nolift}; (I) the Weyl element and the
integrals of Proposition~\ref{prop:notsheaf}; (J)
Proposition~\ref{prop:product}; (K) the Fourier identity of
Proposition~\ref{prop:autoeq}; (L) the least rank $28$ over the monomials. It
prints \texttt{30 checks passed, 0 failed}. The second checks the numerical
statements of \S\ref{ssec:smooth}, among them the moment formula, the
congruence for every $d<2400$, the four rows of \cite[Table~3]{BMB}, the first
candidates of \cite[Table~4]{BMB}, and the bounds on $h^{1}(N_{S/X})$; it
prints \texttt{25 checks passed, 0 failed}.

\subsection*{Data availability}
The two programs and the sources of this note are in the repository above;
the manuscript's code is archived on Zenodo under
\href{https://doi.org/10.5281/zenodo.23038895}{\texttt{10.5281/zenodo.23038895}}
(release v2.0.0; all versions
\href{https://doi.org/10.5281/zenodo.22950275}{\texttt{10.5281/zenodo.22950275}}).
The programs were written and run with the assistance of Claude Code
(Anthropic).

\begingroup\raggedright
\begin{thebibliography}{Mar25b}

\bibitem[BMB]{BMB}
D. Bhattacharjee, P. Mandal and U. Bhattacharya,
\emph{Density of the Algebraic Locus of Weil Classes on Abelian Varieties},
manuscript, 2026; code and sources:
\href{https://doi.org/10.5281/zenodo.23038895}{\texttt{doi:10.5281/zenodo.23038895}}.
Theorem numbers refer to the version of 29 September 2026.

\bibitem[BL04]{BL04}
C. Birkenhake and H. Lange,
\emph{Complex Abelian Varieties}, 2nd ed.,
Grundlehren der mathematischen Wissenschaften \textbf{302},
Springer, Berlin, 2004.
\href{https://doi.org/10.1007/978-3-662-06307-1}{\texttt{doi:10.1007/978-3-662-06307-1}}.

\bibitem[Blo72]{Blo72}
S. Bloch,
\emph{Semi-regularity and de Rham cohomology},
Invent. Math. \textbf{17} (1972), 51--66.
\href{https://doi.org/10.1007/BF01390023}{\texttt{doi:10.1007/BF01390023}}.

\bibitem[Bus19]{Bus19}
N. Buskin,
\emph{Every rational Hodge isometry between two K3 surfaces is algebraic},
J. Reine Angew. Math. \textbf{755} (2019), 127--150.
\href{https://doi.org/10.1515/crelle-2017-0027}{\texttt{doi:10.1515/crelle-2017-0027}}.

\bibitem[Ful98]{Ful98}
W. Fulton,
\emph{Intersection Theory}, 2nd ed.,
Ergebnisse der Mathematik und ihrer Grenzgebiete (3) \textbf{2},
Springer, New York, 1998.
\href{https://doi.org/10.1007/978-1-4612-1700-8}{\texttt{doi:10.1007/978-1-4612-1700-8}}.

\bibitem[Mar25a]{Mar25a}
E. Markman,
\emph{Cycles on abelian $2n$-folds of Weil type from secant sheaves on
abelian $n$-folds}, preprint, arXiv:2502.03415 (2025).
\href{https://doi.org/10.48550/arXiv.2502.03415}{\texttt{doi:10.48550/arXiv.2502.03415}}.

\bibitem[Mar25b]{Mar25b}
E. Markman,
\emph{Secant sheaves and Weil classes on abelian varieties},
preprint, arXiv:2509.23403 (2025).
\href{https://doi.org/10.48550/arXiv.2509.23403}{\texttt{doi:10.48550/arXiv.2509.23403}}.

\bibitem[Mor84]{Mor84}
D. R. Morrison,
\emph{On K3 surfaces with large Picard number},
Invent. Math. \textbf{75} (1984), no.~1, 105--121.
\href{https://doi.org/10.1007/BF01403093}{\texttt{doi:10.1007/BF01403093}}.

\bibitem[Muk81]{Muk81}
S. Mukai,
\emph{Duality between $D(X)$ and $D(\hat{X})$ with its application to
Picard sheaves},
Nagoya Math. J. \textbf{81} (1981), 153--175.
\href{https://doi.org/10.1017/S002776300001922X}{\texttt{doi:10.1017/S002776300001922X}}.

\bibitem[Var23]{Var23}
M. Varesco,
\emph{Hodge similarities, algebraic classes, and Kuga--Satake varieties},
Math. Z. \textbf{305} (2023), paper no.~69.
\href{https://doi.org/10.1007/s00209-023-03390-8}{\texttt{doi:10.1007/s00209-023-03390-8}}.

\end{thebibliography}
\endgroup

\end{document}
